\chapter{Architecture of the 1-bit ALU Slice}
\label{ch:arch}

The whole design rests on one carefully built block: the 1-bit ALU slice.
This chapter develops that slice bottom-up --- full adder, multiplexer,
arithmetic unit, logic unit --- and shows how they compose into the
complete slice, giving the reasoning behind each choice. Every schematic
here is redrawn as vector art; the corresponding Verilog is listed
alongside so the figure and the code can be read together.

\section{Structural decomposition}
The 1-bit slice is partitioned into three cooperating sub-blocks,
selected by the two halves of \sel{}:
\begin{itemize}
  \item an \textbf{arithmetic circuit} that produces the adder result
        $D_i$ and carry-out, controlled by \code{S1\,S0};
  \item a \textbf{logic circuit} that produces the bitwise result $E_i$,
        also controlled by \code{S1\,S0}; and
  \item an output \textbf{4:1 multiplexer} that selects between $D_i$,
        $E_i$, and the two shifted versions of $A$, controlled by
        \code{S3\,S2}.
\end{itemize}
\Cref{fig:alu1top} shows the composition. The key architectural decision
is this two-level select: \code{S3\,S2} picks the \emph{class} of
operation (arithmetic / logic / shift) at the output mux, while
\code{S1\,S0} picks the \emph{specific} operation inside the arithmetic
and logic units. This lets both internal units run in parallel and share
the same lower select bits, and it is exactly what makes the slice
replicable --- no bit position needs special-case control.

\begin{figure}[htbp]
\centering
\begin{tikzpicture}[font=\small]
  % Arithmetic + logic blocks
  \node[arith,minimum width=28mm] (ar) at (0,1.4)
        {1-bit Arithmetic\\ circuit};
  \node[logic,minimum width=28mm] (lo) at (0,-1.4)
        {1-bit Logic\\ circuit};

  % Output mux
  \muxfour{om}{4.4}{2.5}

  % Bundled inputs into the two units
  \draw[bus] ($(ar.west)+(-2.2,0)$)
        node[left,buslbl,align=right]{A$_i$, B$_i$,\\ Cin, S1 S0}
        -- (ar.west);
  \draw[bus] ($(lo.west)+(-2.2,0)$)
        node[left,buslbl,align=right]{A$_i$, B$_i$,\\ S1 S0}
        -- (lo.west);

  % arithmetic -> mux in0, logic -> mux in1
  \draw[sig] (ar.east) -- ($(ar.east)+(0.7,0)$)
        |- node[pos=0.9,above,buslbl]{$D_i$} (om-in0);
  \draw[sig] (lo.east) -- ($(lo.east)+(1.1,0)$)
        |- node[pos=0.9,above,buslbl]{$E_i$} (om-in1);
  % shifts into mux
  \draw[sig] ($(om-in2)+(-1.6,0)$) node[left,buslbl]{shr $A$} -- (om-in2);
  \draw[sig] ($(om-in3)+(-1.6,0)$) node[left,buslbl]{shl $A$} -- (om-in3);
  % select
  \draw[sig] (om-sel) -- ($(om-sel)+(0,-0.6)$) node[below,buslbl]{\code{S3 S2}};

  % output + cout
  \draw[bus] (om-out) -- ($(om-out)+(1.4,0)$) node[right,buslbl]{$F_i$};
  \draw[sig] ($(ar.east)+(0,-0.35)$) -- ($(ar.east)+(2.0,-0.35)$)
        node[right,buslbl]{Cout$_i$};
\end{tikzpicture}
\caption{Top-level datapath of the 1-bit ALU slice. \code{S1\,S0} selects
the function inside the arithmetic and logic units; \code{S3\,S2} selects
the output class (arithmetic $D_i$, logic $E_i$, shift-right or
shift-left of $A$) at the output 4:1 multiplexer.}
\label{fig:alu1top}
\end{figure}

\lstinputlisting[style=vlog,caption={\texttt{alu\_1bit.v} --- top-level
1-bit slice.},label={lst:alu1}]{rtl/alu_1bit.v}

Note in \cref{lst:alu1} that the shift inputs to the output mux
(\code{a\_shr\_1bit}, \code{a\_shl\_1bit}) are derived locally from
\code{a\_i}. In a single bit this is a placeholder for the inter-slice
shift wiring that is realised at the 32-bit level (\cref{ch:part2a}).

\section{The 4:1 multiplexer}
Every selection point in the slice --- inside the arithmetic unit, inside
the logic unit, and at the output --- is a 4:1 multiplexer. A single
parameter-free \code{mux4x1} module is therefore written once and reused,
which is a small but real instance of the reuse principle behind the whole
design.

\begin{figure}[htbp]
\centering
\begin{tikzpicture}[font=\small]
  \muxfour{m}{0}{1.1}
  \draw[sig] ($(m-in0)+(-1.2,0)$) node[left,buslbl]{$a$ (sel=00)} -- (m-in0);
  \draw[sig] ($(m-in1)+(-1.2,0)$) node[left,buslbl]{$b$ (sel=01)} -- (m-in1);
  \draw[sig] ($(m-in2)+(-1.2,0)$) node[left,buslbl]{$c$ (sel=10)} -- (m-in2);
  \draw[sig] ($(m-in3)+(-1.2,0)$) node[left,buslbl]{$d$ (sel=11)} -- (m-in3);
  \draw[bus] (m-out) -- ++(1.2,0) node[right,buslbl]{$y$};
  \draw[sig] (m-sel) -- ++(0,-0.7) node[below,buslbl]{sel[1:0]};
\end{tikzpicture}
\caption{The reusable 4:1 multiplexer. The two select bits route one of
the four data inputs to the output.}
\label{fig:mux}
\end{figure}

\lstinputlisting[style=vlog,caption={\texttt{mux4x1.v} --- reused 4:1
multiplexer.},label={lst:mux}]{rtl/mux4x1.v}

\section{The full adder}
The arithmetic unit is built around a single-bit full adder implementing
\begin{align}
  \text{sum}  &= A \oplus B \oplus C_{in}, \label{eq:fasum}\\
  \text{cout} &= (A\cdot B) + \bigl(C_{in}\cdot(A\oplus B)\bigr).
  \label{eq:facout}
\end{align}
\Cref{fig:fa} is the gate-level realisation of
\cref{eq:fasum,eq:facout}: two XOR gates form the sum, and an
AND-AND-OR network forms the carry, sharing the first XOR ($A\oplus B$)
between the two halves.

\begin{figure}[htbp]
\centering
\begin{circuitikz}[scale=0.9,transform shape,font=\small]
  % Level-1 gates
  \node[xor port]  (x1) at (2.4,3.0) {};
  \node[and port]  (a1) at (2.4,0.6) {};
  % Level-2 gates
  \node[xor port]  (x2) at (5.6,2.6) {};
  \node[and port]  (a2) at (5.6,1.4) {};
  \node[or port]   (o1) at (8.2,1.0) {};

  % Primary inputs
  \node (A) at (0,3.3) {\code{A}};
  \node (B) at (0,2.5) {\code{B}};
  \node (C) at (0,1.0) {\code{Cin}};

  % A,B into x1
  \draw (A) -- (x1.in 1);
  \draw (B) -- (x1.in 2);
  % A,B into a1
  \draw (A) -- ++(0.6,0) |- (a1.in 1);
  \draw (B) -- ++(0.4,0) |- (a1.in 2);
  % x1 -> t1 node, fan to x2.in1 and a2.in1
  \draw (x1.out) -- ++(0.4,0) coordinate (t1);
  \node[circle,fill,inner sep=1pt] at (t1) {};
  \draw (t1) |- (x2.in 1);
  \draw (t1) |- (a2.in 1);
  % Cin -> x2.in2 and a2.in2
  \draw (C) -| ($(x2.in 2)+(-0.6,0)$) -- (x2.in 2);
  \draw (C) -- ++(3.6,0) |- (a2.in 2);
  % sum output
  \draw (x2.out) -- ++(0.8,0) node[right]{\code{sum}\,$=D_i$};
  % carry network
  \draw (a1.out) -- ++(0.5,0) |- (o1.in 1);
  \draw (a2.out) -- (o1.in 2);
  \draw (o1.out) -- ++(0.6,0) node[right]{\code{cout}\,$=$\,Cout$_i$};
\end{circuitikz}
\caption{Gate-level full adder. The shared term $A\oplus B$ feeds both the
sum's second XOR and the carry's second AND, minimising gate count.}
\label{fig:fa}
\end{figure}

\lstinputlisting[style=vlog,caption={\texttt{full\_adder.v}.},
label={lst:fa}]{rtl/full_adder.v}

\section{The arithmetic circuit}
The arithmetic unit turns the single adder into eight operations by
choosing \emph{what the adder's second operand is}. A 4:1 multiplexer
drives the adder's $B$ input with one of four values selected by
\code{S1\,S0}, and \bus{Cin} then distinguishes the two operations within
each opcode (\cref{tab:arithmux}). \Cref{fig:arith} shows the structure.

\begin{table}[htbp]
\centering
\caption{Operand-select encoding of the arithmetic unit. Feeding the adder
$0$, $B$, $\overline{B}$ or $1$ (with the carry-in) realises transfer,
add, subtract and increment/decrement.}
\label{tab:arithmux}
\small
\begin{tabular}{@{}cc l l@{}}
\toprule
\code{S1} & \code{S0} & Adder $B$-input & Resulting function ($C_{in}=0/1$)\\
\midrule
0 & 0 & $0$              & $A$ \,/\, $A+1$ \\
0 & 1 & $B$              & $A+B$ \,/\, $A+B+1$ \\
1 & 0 & $\overline{B}$   & $A+\overline{B}$ \,/\, $A+\overline{B}+1=A-B$ \\
1 & 1 & $1$              & $A-1$ \,/\, $A$ \\
\bottomrule
\end{tabular}
\end{table}

\begin{figure}[htbp]
\centering
\begin{tikzpicture}[font=\small]
  \muxfour{bm}{0}{1.1}
  \draw[sig] ($(bm-in0)+(-1.5,0)$) node[left,buslbl]{$0$} -- (bm-in0);
  \draw[sig] ($(bm-in1)+(-1.5,0)$) node[left,buslbl]{\code{B}} -- (bm-in1);
  \draw[sig] ($(bm-in2)+(-1.5,0)$) node[left,buslbl]{$\overline{\code{B}}$} -- (bm-in2);
  \draw[sig] ($(bm-in3)+(-1.5,0)$) node[left,buslbl]{$1$} -- (bm-in3);
  \draw[sig] (bm-sel) -- ++(0,-0.6) node[below,buslbl]{\code{S1 S0}};

  \node[blk,arith,minimum width=20mm,minimum height=14mm] (fa) at (4.2,0)
        {Full\\ adder};
  \draw[sig] (bm-out) -- ($(fa.west)+(0,0.35)$)
        node[midway,above,buslbl]{$B'$};
  \draw[sig] ($(fa.west)+(-1.3,-0.35)$) node[left,buslbl]{\code{A}}
        -- ($(fa.west)+(0,-0.35)$);
  \draw[sig] ($(fa.south)+(0,-0.9)$) node[below,buslbl]{\code{Cin}}
        -- (fa.south);
  \draw[bus] ($(fa.east)+(0,0.25)$) -- ($(fa.east)+(1.4,0.25)$)
        node[right,buslbl]{$D_i$};
  \draw[sig] ($(fa.east)+(0,-0.35)$) -- ($(fa.east)+(1.4,-0.35)$)
        node[right,buslbl]{Cout$_i$};
\end{tikzpicture}
\caption{Arithmetic circuit: a 4:1 multiplexer selects the adder's second
operand ($0$, $B$, $\overline{B}$ or $1$); the full adder then combines it
with $A$ and the carry-in. Subtraction is two's-complement:
$\overline{B}+C_{in}$.}
\label{fig:arith}
\end{figure}

\lstinputlisting[style=vlog,caption={\texttt{arithmetic\_circuit.v}.},
label={lst:arith}]{rtl/arithmetic_circuit.v}

\paragraph{Design decision --- two's-complement subtraction.}
Rather than build a dedicated subtractor, subtraction reuses the same
adder: $A-B = A + \overline{B} + 1$. Selecting $\overline{B}$ at the mux
and driving $C_{in}=1$ yields $A-B$, while $C_{in}=0$ yields the
subtract-with-borrow result $A+\overline{B}$. This is why the functional
table pairs each opcode with a \bus{Cin} value, and it keeps the slice to
a single adder.

\section{The logic circuit}
The logic unit computes the four bitwise functions of $A$ and $B$ in
parallel and selects one with a 4:1 multiplexer on \code{S1\,S0}
(\cref{fig:logic}).

\begin{figure}[htbp]
\centering
\begin{tikzpicture}[font=\small]
  \node[logic] (g0) at (0,1.6) {\code{A\,\&\,B}};
  \node[logic] (g1) at (0,0.6) {\code{A\,|\,B}};
  \node[logic] (g2) at (0,-0.4){\code{A\,\^{}\,B}};
  \node[logic] (g3) at (0,-1.4){\code{\~{}A}};
  \muxfour{lm}{3.6}{1.1}
  \draw[sig] (g0.east) -- (lm-in0);
  \draw[sig] (g1.east) -- (lm-in1);
  \draw[sig] (g2.east) -- (lm-in2);
  \draw[sig] (g3.east) -- (lm-in3);
  \draw[sig] (lm-sel) -- ++(0,-0.6) node[below,buslbl]{\code{S1 S0}};
  \draw[bus] (lm-out) -- ++(1.3,0) node[right,buslbl]{$E_i$};
  \draw[sig] (-2.2,1.6) node[left,buslbl]{\code{A,B}} -- (g0.west);
  \draw[sig] (-2.2,1.6) -- ++(0.3,0) |- (g1.west);
  \draw[sig] (-2.2,1.6) -- ++(0.15,0) |- (g2.west);
  \draw[sig] (-2.2,1.6) -- ++(0.05,0) |- (g3.west);
\end{tikzpicture}
\caption{Logic circuit: AND, OR, XOR and NOT are computed in parallel and
a 4:1 multiplexer selects the required one.}
\label{fig:logic}
\end{figure}

\lstinputlisting[style=vlog,caption={\texttt{logic\_circuit.v}.},
label={lst:logic}]{rtl/logic_circuit.v}

\section{Why this architecture}
The slice is deliberately regular. All three sub-blocks are combinational,
share the lower select bits, and expose exactly one horizontal carry
signal to the next slice. That regularity is what allows a single verified
cell to be tiled thirty-two times with nothing more than a carry chain and
two shift connections --- the subject of \cref{ch:part2a}. The trade-off,
examined there, is that the resulting carry path is a \emph{ripple} chain,
which becomes the critical path of the 32-bit design.
